\chapter{Fourth Log-Curvature Certificate and Global Second-Order Closure}
\label{chap:g024-global-second-order}

\status{SOH-G024 computer-assisted interval lemma plus exact downstream proof. The first and second complete-monotonicity inequalities for $H_y$ are now global; derivative orders $m\ge3$ remain OPEN. The interval component uses outward \texttt{mpmath.iv} arithmetic and an analytic infinite-theta-tail enclosure. This chapter does not claim RH.}

\section{Purpose}

Chapter~53 proves
\[
 \boxed{H_y''(q)>0\qquad(q\ge1/9,\ 0<|y|<1/2)}.
\]
This chapter closes the remaining compact second-order core. The only computer-assisted ingredient is a fourth-log-curvature inequality for the Riemann kernel; every downstream bridge and Riccati step is exact.

\section{Fourth log-curvature lemma}

Let
\[
 K(t)=\frac12\Phi(|t|),
 \qquad
 L(t)=-\log K(t).
\]
The new certified kernel inequality is
\[
 \boxed{L''''(t)<20L''(t)\qquad(t\in\mathbb R).}
 \tag{54.1}
\]
Since $L$ is even, it is enough to cover $t\ge0$.

For the theta channels
\[
 \phi_n(t)=4a_ne^{5t/2}(2r_n-3)e^{-r_n},
 \qquad
 r_n=a_ne^{2t},\quad a_n=\pi n^2,
\]
write
\[
 \phi_n^{(k)}(t)=4a_ne^{5t/2-r_n}Q_k(r_n).
\]
The exact derivative polynomials through order four are
\[
 Q_0=2r-3,
\]
\[
 Q_1=-4r^2+15r-\frac{15}{2},
\]
\[
 Q_2=8r^3-56r^2+\frac{165}{2}r-\frac{75}{4},
\]
\[
 Q_3=-16r^4+180r^3-529r^2+\frac{1635}{4}r-\frac{375}{8},
\]
and
\[
 Q_4=32r^5-528r^4+2588r^3-4256r^2+\frac{15465}{8}r-\frac{1875}{16}.
\]
If $A_k=\Phi^{(k)}/\Phi$, then
\[
 L''=A_1^2-A_2
\]
and
\[
 L''''=-A_4+4A_3A_1+3A_2^2-12A_2A_1^2+6A_1^4.
\]

\section{Certified compact interval}

The accompanying certificate evaluates the explicit $n=1,\ldots,4$ theta terms with outward interval arithmetic on
\[
 0\le t\le\frac25.
\]
All omitted $n\ge5$ contributions are enclosed analytically. Using $\pi<22/7$, $e^3>20$, and the direct polynomial envelopes, the derivative-order tail bounds are
\[
 4\cdot10^{-28},\quad
 2\cdot10^{-25},\quad
 8\cdot10^{-23},\quad
 5\cdot10^{-20},\quad
 3\cdot10^{-17}.
\]
Adaptive bisection fails closed unless every accepted interval has a strictly positive lower enclosure for
\[
 20L''-L''''.
\]
The emitted receipt records the number of certified boxes, maximum subdivision depth, and smallest outward lower margin.

\section{Analytic tail}

For $t\ge2/5$, write
\[
 \Phi=\phi_1(1+\rho),
 \qquad
 \rho=\sum_{n\ge2}\frac{\phi_n}{\phi_1}.
\]
Because $e^{4/5}>2$ and $\pi>3$, the dominant-channel variable satisfies $r_1>6$. For
\[
 P_1=20(-g_1'')-(-g_1'''')
\]
one has
\[
 P_1-390=
 \frac{1024r^5-12384r^4+52800r^3-106128r^2+92880r-31590}
 {(2r-3)^4}.
\]
After $r=x+6$, the numerator is
\[
 1024x^5+18336x^4+124224x^3+381168x^2+457488x+22842>0,
\]
so
\[
 \boxed{P_1>390}.
\]

The channel derivative bounds for $r\ge6$ give conservative Bell envelopes for the derivatives of $\rho$. Together with the geometric theta suppression, they imply
\[
 |\rho|<10^{-6},\quad
 |\rho'|<\frac1{19000},\quad
 |\rho''|<\frac9{1000},\quad
 |\rho'''|<\frac{17}{10},\quad
 |\rho''''|<378.
\]
Thus, for $h=\log(1+\rho)$,
\[
 |h''|<\frac1{100},
 \qquad
 |h''''|<379,
\]
and
\[
 |h''''-20h''|<380.
\]
Since $L=-g_1-h+\mathrm{constant}$,
\[
 \boxed{20L''-L''''>390-380=10}
\]
for every $t\ge2/5$. This completes the full-line proof of (54.1) at the stated computer-assisted level.

\section{Peano bridge identity}

For the bridge observables
\[
 A_u(r)=L'(u+r)+L'(u-r)
 =\int_{r-u}^{r+u}L''(s)\,ds
\]
and
\[
 B_u(r)=L''(u+r)+L''(u-r),
\]
the trapezoid error is exactly
\[
 \boxed{
 uB_u(r)-A_u(r)
 =\frac12\int_{-u}^{u}(u^2-v^2)L''''(r+v)\,dv.
 }
\]
Using (54.1) and the existing envelope $L''(s)<21e^{2|s|}$,
\[
 \boxed{
 uB_u(r)-A_u(r)
 <280u^3e^{2u}e^{2|r|}.
 }
\]

For the normalized bridge measure,
\[
 R=-\frac{C'}C=\mathbb E[A_u],
 \qquad
 R'=\mathbb E[B_u]-\operatorname{Var}(A_u),
\]
so
\[
 \boxed{uR'-R\le\mathbb E[uB_u-A_u].}
\]

\section{Compact-core Riccati bound}

Put
\[
 a=|y|,
 \qquad
 T=2a\tanh(2au),
 \qquad
 N=R-T,
 \qquad
 S_y(q)=\frac{N(u)}{2u},\quad q=u^2.
\]
Then
\[
 S_y'(q)=\frac{uN'(u)-N(u)}{4u^3}.
\]
Inside the former compact core $u\le1/3$, the sharpened bridge MGF and
\[
 2|r|\le3r^2+\frac13
\]
give
\[
 \mathbb E[e^{2|r|}]
 \le e^{1/3}\left(\frac{17}{14}\right)^{3/2}.
\]
Therefore the bridge part contributes at most
\[
 70e\left(\frac{17}{14}\right)^{3/2}
\]
to $S_y'$. For the tilt, if $x=2au$ then
\[
 \frac{d}{dx}\bigl(\tanh x-x\operatorname{sech}^2x\bigr)
 =2x\operatorname{sech}^2x\tanh x\le2x^2,
\]
so its contribution is at most $1/6$. Hence
\[
 S_y'(q)
 <70e\left(\frac{17}{14}\right)^{3/2}+\frac16.
\]

Using the exact rational bounds
\[
 e<\frac{87}{32},
 \qquad
 \left(\frac{17}{14}\right)^{3/2}<\frac{47}{35},
\]
one obtains
\[
 \boxed{S_y'(q)<\frac{12275}{48}.}
\]
The sharpened first-order theorem gives
\[
 \boxed{S_y(q)>\frac{33}{2}},
\]
therefore
\[
 S_y(q)^2>\frac{1089}{4}.
\]
The strict gap is
\[
 \boxed{
 \frac{1089}{4}-\frac{12275}{48}
 =\frac{793}{48}>0.
 }
\]
Consequently
\[
 \boxed{
 \frac{H_y''(q)}{H_y(q)}
 =S_y(q)^2-S_y'(q)>0
 }
\]
for $0<q\le1/9$, and $q=0$ follows by continuity.

\section{SOH-G024 global second-order theorem}

Combining the compact-core result with Chapter~53 gives:

\begin{theorem}[SOH-G024 global second-order theorem]
For every external Dimitrov--Xu tilt with $0<|y|<1/2$,
\[
 \boxed{H_y''(q)>0\qquad(q\ge0).}
\]
The theorem depends on the declared computer-assisted interval certificate for (54.1).
\end{theorem}

Thus the first two non-trivial complete-monotonicity inequalities are global:
\[
 \boxed{H_y'(q)<0},
 \qquad
 \boxed{H_y''(q)>0}.
\]

\gap{Derivative orders $m\ge3$ remain OPEN. Therefore full complete monotonicity, strict external Fourier positivity, Wiener density, SOH-G003, SOH-C005, PF$_3$, PF$_\infty$, and RH remain OPEN.}
